% ================= 8. CONCLUSIONS & NEXT STEPS =================
\section{Conclusions \& Next Steps}
\label{sec:conclusions}

\subsection{Conclusions}
\label{sec:conclusions_summary}

The sizing closure of Section~\ref{sec:sizing} has converged to a fixed point:
\texttt{analysis/SIZING.py} reports a cube mass of \qty{1.772}{\kilo\gram} at
which the mass budget, the jump-up inertia requirement and the ballasted
reaction-wheel design agree to within \qty{0.1}{\gram}, with the wheel
delivering \qty{111.7}{\percent} of the corner-case target and a trim
authority that absorbs a heavier-than-predicted structure comfortably (Section
\ref{sec:sizing_verification}). Both prototypes named in the methodology of
Section~\ref{sec:methodology} exist as concept articles: the single-face 1-DoF
rig is built to its released drawing (Figure~\ref{fig:rig_2d}), and a first
printed concept of the three-axis cube frame demonstrates the two-part
architecture and internal packaging (Figure~\ref{fig:cube_3d}), though its
final dimensions remain open pending the packaging study below. The dynamics
are derived for both the 2D and 3D cases, the LQR baseline controller and its
executable logic are stated in full (Section~\ref{sec:control_algorithms}),
and the electrical architecture --- power tree, bus topology, schematic sheet
set --- is fixed at the level needed to build a harness (Section
\ref{sec:electrical}).

What remains open is concentrated in a small number of places rather than
spread evenly across the design. The structural mass in
Table~\ref{tab:massbudget} is still a \qty{34}{\percent} allowance rather than
a measured figure, which is the dominant uncertainty in the sizing closure.
The final motor and the three-dimensional frame's driving dimensions
(Table~\ref{tab:cad_params}) are not yet confirmed. Appendix
\ref{app:control_algorithms}'s gains, timing budget and several guard
conditions are stated as TBD pending plant identification and the simulation
work of Section~\ref{sec:sim_environment}. And the hardware validation of
Section~\ref{sec:validation_results} below is, at the time of writing, a set
of defined criteria rather than a set of results.

\subsection{Validation Results to Date}
\label{sec:validation_results}

This subsection reports evidence against the validation stages of
Section~\ref{sec:test_strategy} and the acceptance criteria of
Table~\ref{tab:traceability}, in the order the stages are entered. Where no
result exists yet, the criterion is stated so the gap is explicit rather than
silent.

\subsubsection{1-DoF Control Validation}
\label{sec:test_1dof}

The single-axis rig is the protected experimental item of
Section~\ref{sec:schedule}, and its balance result is the evidence carried into
this document. Gate G3 requires the rig to balance for at least 60~s unassisted.

Beyond the duration criterion, three quantities are recorded because they
determine whether the result transfers to the integrated cube rather than merely
holding on the rig. Recovery from a defined impulse establishes that the balance
is a stabilised equilibrium and not a fortunate initial condition. Steady-state
wheel-speed drift measures momentum accumulation: a controller that balances
while monotonically spinning up is one that will saturate, and the drift rate
predicts when. Control effort is compared against the continuous and peak current
limits of the driver given in Appendix~\ref{app:bom}, since a balance sustained
only by exceeding the continuous rating is not a balance that survives to the
demonstration.

\subsubsection{Edge, Corner and Jump-Up Validation}
\label{sec:test_milestones}

The three hardware milestones follow the objective ordering of
Section~\ref{sec:objectives}, each entered only on the passing of the previous.

\paragraph{Edge balance.}
\label{sec:test_edge}
Regulation about a
line-contact equilibrium, and the hard-kill gate for the jump-up objective. G6
requires balance for at least 60~s with disturbance recovery. Recorded:
sustained duration; the maximum initial tilt from which the controller
recovers, which is the experimental region of attraction and is compared
against the simulated figure from Section~\ref{sec:sim_environment}; the
disturbance-rejection bandwidth; and wheel-speed drift as in
Section~\ref{sec:test_1dof}. Numeric targets beyond the gate criterion are
\textcolor{red}{TBD}.

\paragraph{Corner balance and commanded slew.}
\label{sec:test_corner}
Regulation about a point-contact equilibrium on all three axes simultaneously,
which is the case the wheel was sized against and therefore the first test in
which the sizing assumption of Section~\ref{sec:sizing_verification} is
exercised directly. Recorded: sustained duration, per-axis control effort, and
cross-axis coupling, the last being absent from the 2D model and observable
only here.
\label{sec:test_slew}
Commanded rotation between equilibria is verified in the same stage: the cube is
commanded through a slew and the tracking error against the reference recorded,
together with the settling behaviour on arrival. Targets \textcolor{red}{TBD}.

\paragraph{Jump-up.}
\label{sec:test_jumpup}
The face-to-edge transition,
captured on video at a frame rate sufficient to resolve the manoeuvre, per G8.
This test closes the loop opened in Section~\ref{sec:sizing_verification}. The
wheel speed at brake release is measured and compared against the value
predicted before fabrication, and the two routes of Section
\ref{sec:sizing_verification} are re-evaluated against the measured brake
torque. The discrepancy is the quantitative measure of the inelastic-collision
idealisation, and its measurement is the experimental refinement of $\omega_w$
that Section~\ref{sec:jumpup} anticipates. Recorded: wheel speed at release,
achieved versus predicted; whether the edge is reached and held; and, on a
successful transition, the time to re-stabilise, which is the quantity that
determines whether chained jump-up --- the walking objective --- is reachable.

\subsubsection{Control Performance Evaluation}
\label{sec:test_perf}

The metrics used above are defined once here so that simulated and measured
results are expressed in the same terms and remain comparable across stages.

\begin{table}[htbp]
\centering
\caption[Control performance metrics]{Control performance metrics. Definitions and methods are fixed;
targets follow from the closed design point.}
\label{tab:perf_metrics}
\small
\setlength{\tabcolsep}{4pt}
\begin{tabular}{L{2.8cm} L{4.7cm} L{4.1cm} c}
\toprule
Metric & Definition & Measurement method & Target \\
\midrule
Recovery angle & Largest initial tilt from which the equilibrium is regained & Ramped initial-condition release & \textcolor{red}{TBD} \\
Settling time & Time to enter and remain within a stated band about the equilibrium & Step or impulse response log & \textcolor{red}{TBD} \\
Disturbance-rejection bandwidth & Frequency to which body disturbances are attenuated & Swept or impulsive disturbance & \textcolor{red}{TBD} \\
Slew tracking error & Peak and RMS deviation from the commanded trajectory & Reference versus estimated attitude & \textcolor{red}{TBD} \\
Control effort & RMS and peak wheel current over the run & Driver current telemetry & Within Appendix~\ref{app:bom} limits \\
Momentum drift & Rate of change of mean wheel speed while regulating & Encoder log over a sustained run & \textcolor{red}{TBD} \\
\bottomrule
\end{tabular}
\end{table}

\subsubsection{Milestone Acceptance Criteria}
\label{sec:test_acceptance}

Table~\ref{tab:acceptance} states the condition under which each objective of
Table~\ref{tab:traceability} is accepted as demonstrated, and the subsection that
produces the evidence. It is the closing element of the traceability chain that
runs from the objectives of Section~\ref{sec:objectives}, through the required
capabilities of Section~\ref{sec:requirements}, to the test that confirms them.

\begin{table}[htbp]
\centering
\caption[Objective acceptance criteria and supporting evidence]{Objective acceptance criteria and the evidence that supports each.}
\label{tab:acceptance}
\small
\setlength{\tabcolsep}{4pt}
\begin{tabular}{L{3.0cm} L{6.4cm} c c}
\toprule
Objective & Acceptance condition & Evidence & Gate \\
\midrule
Edge balance & Balanced $\geq$60 s unassisted with recovery from a defined disturbance & Sec.~\ref{sec:test_edge} & G6 \\
Corner balance & Point-contact regulation sustained on all three axes & Sec.~\ref{sec:test_corner} & G7 \\
Controlled rotation & Commanded slew completed within the stated tracking error & Sec.~\ref{sec:test_slew} & G7 \\
Jump-up & Face-to-edge transition achieved and captured on video & Sec.~\ref{sec:test_jumpup} & G8 \\
Walking (stretch) & Chained jump-up with re-stabilisation between transitions & Sec.~\ref{sec:test_jumpup} & --- \\
\bottomrule
\end{tabular}
\end{table}

\subsection{Immediate Next Steps}
\label{sec:next_steps}

In priority order, the work that closes the open items of
Section~\ref{sec:conclusions_summary}:

\begin{enumerate}
  \item \textbf{Close the packaging study} (Table~\ref{tab:cad_params}): the
        minimum internal clear volume $V_{\min}$, the motor axis offsets and
        the inter-module clearances are the head of the sizing dependency
        chain $V_{\min}\to L\to M\to h_w\to I_{w,\text{target}}$, and every
        other open item downstream of it moves again once it closes.
  \item \textbf{Build and run the Simulink/Simscape environment}
        (Section~\ref{sec:sim_environment}) against the CAD of both
        prototypes, and use it to re-derive the gains of
        Table~\ref{tab:gains} once the 1-DoF plant identification of
        Section~\ref{sec:paramid} reports measured parameters.
  \item \textbf{Complete Appendix~\ref{app:control_algorithms}}: the guard
        conditions of Table~\ref{tab:state_transitions}, the gains of
        Table~\ref{tab:gains} and the timing budget of
        Table~\ref{tab:timing_budget} close once steps~1--2 above do.
  \item \textbf{Run the electrical and calibration checks of
        Section~\ref{sec:verification}} in order --- CAN-bus termination and
        rail voltages, then IMU and moteus calibration, then the gradual
        control-authority ramp --- before any unassisted balancing attempt.
  \item \textbf{Populate the decidable-now tables}: fabrication routes and
        print settings (Appendix~\ref{app:cad}), the physical pin numbers of
        the per-sheet tables in Appendix~\ref{app:electrical} once the board
        layout of Section~\ref{sec:layout} is drawn, and the connector/wire
        schedules.
  \item \textbf{Carry the milestone validations of
        Section~\ref{sec:validation_results} to completion} against the gate
        schedule of Table~\ref{tab:gates}, reporting results in place of the
        stated criteria as each stage closes.
\end{enumerate}

The longer-horizon research directions below are explicitly not part of this
punch-list: none of them is a precondition for the milestones above, and the
baseline LQR controller is not contingent on any of them.

% \subsection{NMPC \& Warm-Start Control (Future Implementation)}

% \subsection{Autonomous Stand-up \& Disturbance Rejection}

\subsection{AI-Driven Control \& Dynamic Coefficient Tuning}
\label{sec:ai_control}

\subsubsection{Motivation: Limitations of a Fixed Gain}

The balancing controller of Section~\ref{sec:control2d} computes a single
gain $K_d$ offline, from an $(A,B)$ pair built from the parameters of
Tables~\ref{tab:2d_vars} and \ref{tab:3d_vars}. That gain is optimal for
exactly one plant: the one those numbers describe. It is retained here as the
baseline precisely because it is simple and verifiable, but its optimality is
conditional in three ways that the hardware of this project is expected to
violate.

\begin{itemize}
    \item \textbf{Parametric drift.} $C_b$ and $C_w$ are not obtained from
    CAD and are carried as estimates until the identification tests of
    Table~\ref{tab:params}.
    Bearing friction,
    wiring drag and motor cogging also change with temperature, wear and
    reassembly, so the identified values are a snapshot rather than a
    constant. The inertia tensor $I$ and the centre-of-mass offset $r_{cm}$
    are subject to the same objection at the few-percent level, and $r_{cm}$
    in particular shifts whenever the pack or harness is disturbed.
    \item \textbf{Payload and configuration changes.} Any added mass changes
    $m$, $r_{cm}$ and $I$ simultaneously, and therefore changes $A$. A gain
    tuned for one configuration is not merely suboptimal for another; near a
    fully unstable equilibrium it may be destabilising.
    \item \textbf{Contact and surface variation.} The pivot is modelled as an
    ideal frictionless point or line. A real corner resting on a real surface
    slips, deforms and dissipates energy in a way that is neither in the model
    nor constant across surfaces, and this is the dominant unmodelled effect
    for the walking objective of Section~\ref{sec:objectives}.
\end{itemize}

The standard remedy within the existing design is gain scheduling on
$\hat\theta_b$ (Section~\ref{sec:control2d}), but scheduling interpolates
between gains that were themselves computed from the same assumed parameters.
It addresses the weakening of the small-angle linearisation; it does not
address a plant that differs from the model. The proposal here is to keep the
verified feedback law and adapt its coefficients online.

\subsubsection{Proposed Hybrid Architecture}

The intended structure is deliberately conservative: a classical inner loop
whose stability properties are understood, wrapped by a slow outer loop that
adjusts the inner loop's coefficients. Writing $\phi$ for the tunable
coefficient vector,
\begin{equation}
u[k] = -K_d(\phi[k])\,\hat x[k],
\qquad
\phi[k+1] = \phi[k] + \Delta\phi\big(\hat x[k],\, \hat x[k-1], \dots\big),
\label{eq:hybrid_arch}
\end{equation}
where $\hat x$ is the estimated state of Section~\ref{sec:control2d} and
$\Delta\phi$ is produced by the tuner. Two properties of
\eqref{eq:hybrid_arch} matter more than the choice of learning algorithm.
First, the inner loop remains a linear state feedback at every instant, so
the controller that runs on the hardware is the one already implemented and
tested. Second, $\phi$ is updated on a timescale far slower than the control
period $\Delta t$, which keeps the frozen-gain analysis approximately valid
and prevents the tuner from acting as an unmodelled fast feedback path. Three
candidate tuners are of interest, and they are not equivalent in cost:

\begin{itemize}
    \item \textbf{Bayesian optimisation} over a low-dimensional
    parameterization of $Q$ and $R$, minimising a measured cost from short
    balancing trials \cite{shahriari2016bayesopt}. This is the cheapest option
    and the one with a direct precedent on this class of hardware: automatic
    LQR tuning has been demonstrated on an inverted pendulum with on the order
    of tens of experiments \cite{marco2016automaticlqr}, and a self-tuning LQR
    has been run on the Cubli's own ancestor platform at ETH
    \cite{trimpe2014learningcubli}. It tunes between trials, not within them.
    \item \textbf{Adaptive dynamic programming}, which solves the LQR problem
    online from measured data rather than from $(A,B)$, converging toward the
    optimal gain for the true plant without an explicit model
    \cite{lewis2012adp}. This is the option with the strongest theoretical
    guarantees and the strictest excitation requirements.
    \item \textbf{Reinforcement learning}, in which a policy trained in
    simulation outputs $\Delta\phi$ continuously from the state-error history.
    This is the most capable and the least verifiable, and is treated below as
    the exploratory branch rather than the default.
\end{itemize}

\paragraph{Deployment reservations.} The literature supporting learned control
is overwhelmingly validated in simulation or on well-instrumented hardware,
and the failure mode of an online tuner on an open-loop-unstable plant is a
fall, not a degraded score. Any deployment on this cube should therefore be
gated behind: a projection of $\phi$ onto a set for which the frozen-gain
closed loop is verified stable across the parameter uncertainty; a fallback to
the fixed $K_d$ of Section~\ref{sec:control2d} on any estimator fault or
divergence; and bench trials with the cube tethered. This is future work in
the literal sense --- none of it is a committed deliverable of the present
design, and the baseline controller is not contingent on any of it.

\subsubsection{Module 1: Dynamic Gain Tuning}

The tuner does not emit motor currents. It emits the coefficients from which
$K_d$ is computed, which bounds what a bad update can do. Two
parameterizations are worth comparing:

\begin{itemize}
    \item \textbf{Weight-space tuning.} The policy outputs the diagonal
    entries of the cost matrices,
    $\phi = (\log q_1,\dots,\log q_9,\ \log r_1, \log r_2, \log r_3)$ with
    $Q = \mathrm{diag}(q_i) \succ 0$ and $R = \mathrm{diag}(r_i) \succ 0$, and
    $K_d$ follows from the discrete algebraic Riccati equation. The
    logarithmic parameterization enforces positive definiteness by
    construction, so every $\phi$ the policy can emit corresponds to a
    well-posed LQR problem and, for a stabilisable $(A_d,B_d)$, to a
    stabilising gain for the \emph{modelled} plant. The cost is that a Riccati
    solve must run onboard whenever $\phi$ changes, which is what confines
    updates to a slow outer loop.
    \item \textbf{Gain-space tuning.} The policy perturbs the entries of $K_d$
    directly, or equivalently the $K_p, K_d$ pair of a PD form, avoiding the
    Riccati solve at the price of losing the structural guarantee above. This
    is the practical choice if the outer loop must run fast, and it makes the
    projection step of the previous paragraph mandatory rather than merely
    prudent.
\end{itemize}

For the 3D corner-balancing case the state is
$x = (\theta, \omega_b, \omega_w) \in \mathbb{R}^9$ and $K_d$ is $3\times 9$
(Section~\ref{sec:3dmodel}), so tuning the $54$ free entries directly is not
sensible on the available data. Weight-space tuning reduces the problem to
$12$ coefficients, and symmetry of the three wheel axes reduces it further if
the wheels are treated as identical. A reward or cost of the form
\begin{equation}
c[k] = \hat x^{\mathrm T}[k]\,W_x\,\hat x[k] + \rho\,\|u[k]\|^2 + \sigma\,\|\omega_w[k]\|^2
\label{eq:tuner_cost}
\end{equation}
adds an explicit penalty on stored wheel momentum, which the nominal LQR cost
$J(u)$ of Section~\ref{sec:control2d} does not distinguish from any other
state deviation. This term is the one directly relevant to the saturation
problem raised in Section~\ref{sec:application}: it expresses a preference for
holding attitude with the wheels near zero speed, leaving momentum headroom
for disturbance rejection.

\subsubsection{Module 2: Sim-to-Real Transfer via Domain Randomization}

A policy trained against a single nominal parameter set learns to exploit that
parameter set. Domain randomization instead samples the plant parameters from
a distribution at the start of each training episode
\cite{tobin2017domainrand}, so that the policy is optimised against a family
of plants rather than one, and the transfer method has a substantial record on
physical legged hardware \cite{hwangbo2019anymal}. For this cube the
randomized quantities follow directly from the identification gaps already
documented:

\begin{itemize}
    \item \textbf{Damping:} $C_b, C_w$ over a wide interval, since these are
    estimated rather than measured until the identification of
    Section~\ref{sec:paramid} closes.
    \item \textbf{Inertia and mass distribution:} $I$, $m$ and $r_{cm}$
    perturbed by a few percent, covering print-density variation, fastener
    placement and harness routing.
    \item \textbf{Actuation:} $K_m$ within motor tolerance, plus a current
    saturation limit and a commanded-torque delay, so the policy cannot assume
    an ideal actuator.
    \item \textbf{Contact:} pivot friction and restitution, and a small random
    ground slope --- the term standing in for the surface variation that
    motivates this module.
    \item \textbf{Sensing:} gyro bias walk, accelerometer noise and an
    injected vibration component, matching the estimator degradation
    identified in Section~\ref{sec:background} as the practical limit on
    balancing.
\end{itemize}

The sensing term requires particular care. The accelerometer corruption on this
platform is generated by the controller's own wheel activity, so it is
correlated with the policy's actions rather than independent of them. A
randomization scheme that injects only white sensor noise will understate the
difficulty, and a policy that transfers under it should not be assumed to
transfer on hardware. Whether the randomized distribution covers the real cube
is an empirical question, settled by bench testing; the principal failure mode
of this module is a distribution selected for convenience rather than
representativeness.

\subsubsection{Module 3: Jump-Up Energy Optimisation}

The jump-up analysis of Section~\ref{sec:jumpup_target} treats the manoeuvre as
a spin-up to $\omega_{\max}$ followed by an abrupt stop, with all deviation
from the ideal absorbed into a single brake-amplification factor
$\beta = \tau_b/(\tau_b-\tau_g)$. That model is adequate for sizing and is
deliberately conservative, but it is a poor description of what the mechanism
does: the servo-driven barrier has finite travel, the collision is not
instantaneous, and the resulting torque profile is neither a step nor known in
advance.

The proposal is to treat the torque command over the manoeuvre window as a
profile to be optimised rather than a constant. Parameterizing the
pre-brake and post-impact torque commands by a small vector $\psi$ --- for
instance the timing of brake release relative to the wheel-speed trajectory,
and the shape of the recovery torque applied by the remaining wheels
immediately after impact --- the objective is
\begin{equation}
\min_{\psi}\ \ \mathbb{E}\Big[\, \|\theta(t_f) - \theta^\star\|^2 + \mu\,\|\omega_b(t_f)\|^2 \,\Big]
\quad \text{subject to} \quad
\omega_w \le \omega_{\max},\ \ |u_i| \le u_{\max},
\label{eq:jump_opt}
\end{equation}
where $t_f$ is the instant the controller hands over to the balancing law at
$|\hat\theta_b| < \theta_{\text{th}}$. Minimising residual body rate at
handover is the quantity of interest, because a jump that reaches the
equilibrium with excess rate demands recovery authority the balancing
controller may not have. Because each evaluation of \eqref{eq:jump_opt} is a
physical jump, the sample budget is small and the appropriate method is
Bayesian optimisation rather than a policy-gradient method, with the initial
design taken from the analytical $h_w$ of Equation~\eqref{eq:hw_ideal}.

Two constraints on this module must be stated. First, it presupposes the
measurement described in Section~\ref{sec:jumpup_target}: until the spin-down
test reports both the impulse-equivalent mean
$\bar\tau_b = h_w/\Delta t$ and the peak torque, there is no validated model to
optimise against and no way to distinguish an improved profile from a
mis-measured one. Second, the structural caveat carries over unchanged ---
optimisation must not be permitted to search toward shorter stopping times,
since $\tau_b$ is simultaneously a load on the barrier, bolt, bracket and
spokes, and the sizing value was chosen for the dynamics and not for that load
path. The search space must be bounded by the structural case, not by the
dynamic one.

\subsubsection{Relevance Beyond the Bench}

Two application claims follow, and both are narrower than the general
enthusiasm for learned control would suggest.

For \textbf{CubeSat attitude control}, the relevant fact is that a
spacecraft's inertia tensor is known imprecisely at integration and changes in
flight --- propellant depletion, deployable appendages, thermal distortion ---
while reaction-wheel friction rises with bearing age. These are the orbital
analogues of the parametric drift of Section~\ref{sec:control2d}, and they are
conventionally handled by conservative fixed gains with margin, at a cost in
pointing performance. A tuner that adapts $Q$ and $R$ against measured
performance addresses exactly this gap. The qualification barrier for flying
learned control is high and unaddressed here; the architecture of
\eqref{eq:hybrid_arch} is chosen partly because a bounded outer loop over a
classical inner loop is the form most likely to survive that scrutiny.

For \textbf{low-gravity surface mobility}, the case is stronger, because the
plant is genuinely unknown before arrival. Regolith mechanical properties on a
small body are not characterised to the precision a hopping controller would
want, and the flight-proven hoppers cited in Section~\ref{sec:application}
\cite{yoshimitsu1999minerva, ho2017mascot} move ballistically rather than
under closed-loop control for this reason. The contact randomization of
Module~2 and the profile optimisation of Module~3 map onto that problem
directly: a hop whose landing attitude is controlled, on a surface whose
properties were sampled from a distribution rather than assumed, is the
capability that separates directional walking from a ballistic tumble
\cite{spiridonov2024spacehopper}. The cube tests this under $1$~g, where the
destabilising torque is larger and never relents --- a harder case in that one
respect, as argued in Section~\ref{sec:application}, though it does not
reproduce the low-gravity contact dynamics that would decide the real
manoeuvre.

% \subsection{Single-Wheel Variant / Extensions}

% A reduced-actuation variant balancing the 3D cube with a single reaction wheel,
% following the ETH one-wheel Cubli \cite{hofer2023onewheelcubli}.
